% Einheit 3 von 5 – Zusammengesetzte Flächen, Maßstab und Volumen (bank/flaecheninhalt-durch-integration/e3.jsonl, zusammenbau v0.1)
%% TODO Zeitmarke Einheit 3: Marken-Zeile nicht lesbar
%% TODO Zweigzeile Teil 1: was hier gelernt wird (Satz in Schülersprache aus dem Einheitstitel) – Einheit 3
\einheitenkopf[e3]{Einheit 3 von 5 · Zusammengesetzte Flächen, Maßstab und Volumen}
\zweigzeile{Abitur GK}
\setcounter{aufgabe}{12}

%% TODO Titel: Ich-kann-Satz fehlt in der Bank (Platzhalter: Kettenname) – Nr. 13
%% TODO Anweisung: ein Satz über der ersten Teilaufgabe fehlt in der Bank – Nr. 13
\begin{aufgabe}{Zusammensetzen, Maßstab, Volumen}
%% TODO \rechenplatz{n}: Schrittzahl des Grundfalls fehlt in der Bank (nur bei fester Schreibform, 2.3 b)
\begin{teile}
\teil Teilstück unter dem Bogen von $f(x) = 0{,}5x^2 + 1$ zwischen $x = 0$ und $x = 2$ – Integral oder Formelfläche? \kreuz{Integral}\\ \kreuz{Formelfläche}
\teil Eine Fläche wird oben vom Graphen von $f(x) = x^2$ ($0 \le x \le 2$) und der Strecke von $(2|4)$ nach $(4|0)$ begrenzt, unten von der x-Achse – Inhalt?
\teil Der obere Rand ist $f(x) = 0{,}5x + 1$ für $0 \le x \le 2$ und $g(x) = -0{,}5x^2 + 2x$ für $2 \le x \le 4$; unten die x-Achse, links die y-Achse – Inhalt?
\teil Ein Grundstück liegt zwischen dem Graphen von $f(x) = -0{,}5x^2 + 4x$ und der x-Achse; eine Längeneinheit entspricht $20$ m – Fläche in Hektar?
\teil Eine Betonschwelle hat den Querschnitt zwischen dem Graphen von $f(x) = -0{,}25x^2 + 1$ und der x-Achse (Angaben in dm) und ist $30$ dm lang; ein Kubikdezimeter Beton wiegt $2{,}4$ kg – Masse?
\teil Eine Spielfläche besteht aus der Fläche zwischen dem Graphen von $f(x) = -0{,}25x^2 + 1$ und der x-Achse und einem Halbkreis mit Radius $2$ unter der x-Achse – Inhalt? Runde auf zwei Nachkommastellen.
\teil $f(x) = -x^2 + x + 2$. Jemand rechnet: (1) $2x + 2 = f(x)$ liefert $-1$ und $0$; (2) das Integral von $-1$ bis $0$ über $f(x) - (2x + 2)$ ergibt $\frac{1}{6}$; (3) $0{,}5 \cdot f(0) = 1$; (4) die Zahlen werden ins Verhältnis gesetzt. Deute jeden Schritt geometrisch. \hfill (Abitur 2024 GK)
\end{teile}
\end{aufgabe}

%% TODO Titel: Ich-kann-Satz fehlt in der Bank (Platzhalter: Kettenname) – Nr. 14
%% TODO Anweisung: ein Satz über der ersten Teilaufgabe fehlt in der Bank – Nr. 14
\begin{aufgabe}{Zusammensetzen, Maßstab, Volumen – Fehler finden, Begründen, Anwendung, Darstellungswechsel}
\begin{teile}
\teil Eine Fläche hat die Maßzahl $5{,}5$, eine Längeneinheit entspricht $4$ m – wo ist der Fehler? \rechnung{A &= 5{,}5 \cdot 4 = 22\,\text{m}^2}
\teil Warum ist der Flächenmaßstab das Quadrat des Längenmaßstabs?
\teil Eine Regenrinne hat den Querschnitt zwischen der x-Achse und dem Graphen von $f(x) = 0{,}2x^2 - 0{,}8$ (Angaben in dm) und ist $8$ m lang. Wie viele Liter fasst sie randvoll?
\teil Im Bild ist der Graph von $f(x) = 0{,}5x^2 + 1$. Markiere die Fläche, deren Inhalt das Integral von $0$ bis $2$ über $f(x)$ plus $2 \cdot 3$ angibt.

\begin{ksys}[xmin=-1,xmax=5,ymin=-1,ymax=5,ablesen]\funktion{0.5*\x^2+1}{f}\end{ksys}
\end{teile}
\end{aufgabe}
